\documentclass[11pt]{article}
\usepackage[margin=1in]{geometry}
\usepackage[T1]{fontenc}
\usepackage{lmodern}
\usepackage{amsmath,amssymb}
\usepackage{booktabs}
\usepackage{array}
\usepackage{longtable}
\usepackage[table]{xcolor}
\usepackage{hyperref}
\usepackage{microtype}

\hypersetup{
  colorlinks=true,
  linkcolor=blue!55!black,
  urlcolor=blue!55!black
}

\definecolor{mistakered}{RGB}{185,20,20}
\newcommand{\wrong}[1]{\textcolor{mistakered}{#1}}
\newenvironment{mistake}
  {\begin{quote}\color{mistakered}\textbf{Mistake.}\ }
  {\end{quote}}
\newenvironment{statusnote}
  {\begin{quote}\color{blue!55!black}\textbf{Verified status.}\ }
  {\end{quote}}

\title{Revised Audit Report on Q1\\
The Double-Octic Model, Its Singular Locus, and the Unfinished Resolution}
\author{Prepared from the Q1 polynomial, exact Singular computations,\\
the Rust blowup audit, and \texttt{q1\_expository\_report.pdf}}
\date{July 15, 2026}

\begin{document}
\maketitle

\begin{abstract}
This report reconciles the newly supplied expository PDF with exact calculations for the
given branch octic.  The PDF and the Rust input define exactly the same polynomial
$\Delta=q_4^2-4q_3q_5$.  The PDF correctly identifies seven singular lines and an
isolated point.  It does not, however, correctly compute five of the seven generic
transverse $cA_n$ types.  More seriously, the isolated point is already an ordinary
double point of the double cover, is disjoint from every center in the advertised
divisorial ledger, and lies on none of the three advertised small-resolution divisors.
Consequently the claimed smooth projective model is not obtained by the construction as
written.  Claims found to be false or unsupported are printed in red.
\end{abstract}

\tableofcontents

\section{Executive conclusion}

Let $Y_0$ be the double cover
\[
  Y_0:\quad R^2=\Delta(a,b,c,d)
  \subset \mathbb P(1,1,1,1,4).
\]
The following facts have been checked exactly over $\mathbb Q$.

\begin{enumerate}
  \item The polynomial supplied to the Rust program is exactly the discriminant
  $q_4^2-4q_3q_5$ displayed in the expository PDF.
  \item The branch singular locus contains the seven lines $L_1,\ldots,L_7$ listed in
  the PDF and the isolated point
  \[
    P_*=[4:-2:3:1].
  \]
  \item The generic transverse types in line order are
  \[
    (cA_3,cA_3,cA_3,cA_4,cA_1,cA_5,cA_1),
  \]
  not the sequence stated in the PDF.
  \item The branch Hessian at $P_*$ has nonzero determinant, so $Y_0$ has an ODP above
  $P_*$.  None of the listed centers meets $P_*$, and none of the three listed Weil
  divisors contains the point above $P_*$.
\end{enumerate}

\begin{mistake}
The assertion in Sections 5.4--5.5 of \texttt{q1\_expository\_report.pdf} that the
advertised 36 divisorial blowups followed by the three listed Weil-divisor blowups produce
a smooth threefold is false as written: the ODP above $P_*$ is not resolved.
\end{mistake}

It follows that the correction polynomial $58p^2+82p$, the claimed Picard rank, Euler
characteristic, rigidity, and modularity conclusion cannot be treated as proved by the
provided ledger.  They may ultimately be true for a corrected resolution, but a corrected
global construction and new exceptional-stratum count are required.

\section{Identity of the model}

The PDF defines
\[
q_3=ad(b+c-d)
\]
and homogeneous polynomials $q_4,q_5$ of degrees four and five, respectively, then sets
\[
  \Delta=q_4^2-4q_3q_5.
\]
Expanding this expression in $a,b,c,d$ gives exactly the polynomial in
\texttt{q1\_branch\_polynomial.txt} after the identification
\[
  (a,b,c,d)=(x_0,x_1,x_2,x_3).
\]
Thus none of the discrepancies below can be attributed to a change of coordinates or to a
different octic.

The projective singular components reported by the PDF are
\begin{align*}
L_1&:(a,b), & L_2&:(a-b,c), & L_3&:(a,d),\\
L_4&:(b,c-d), & L_5&:(a,c-d), & L_6&:(c,d),\\
L_7&:(a-d,b+c-d),
\end{align*}
together with
\[
  P_*:(c-3d,b+2d,a-4d).
\]
Direct substitution verifies that $\Delta$ and all four homogeneous partial derivatives
vanish on each listed component.

\begin{statusnote}
The eight-component singular-locus decomposition in Section 4.1 of the PDF is correct.
The earlier Rust report was wrong to present the seven lines as the complete initial
singular locus; its projective linear fast path returned before adding the isolated point.
\end{statusnote}

\section{Exact transverse singularity types}

\subsection{Method}

At a generic rational point of each line, choose an affine transverse plane with coordinates
$u,v$.  Let $f_i(u,v)$ be the restriction of $\Delta$ to that plane.  The Milnor number is
computed in the local ring by
\[
  \mu_i=\dim_{\mathbb Q}\frac{\mathbb Q[[u,v]]}
  {(\partial f_i/\partial u,\partial f_i/\partial v)}.
\]
For these rank-one quadratic plane germs, $\mu_i=n$ is precisely the $A_n$ index, so the
double-cover family has transverse type $cA_n$.

The exact slices used by \texttt{q1\_exact\_checks.sing} are shown below.

\begin{center}
\small
\begin{tabular}{@{}cll@{}}
\toprule
Line & Generic point & Transverse substitution \\
\midrule
$L_1$ & $[0:0:1:-1]$ & $(a,b,c,d)=(u,v,1,-1)$ \\
$L_2$ & $[1:1:0:2]$ & $(a,b,c,d)=(1,1+u,v,2)$ \\
$L_3$ & $[0:1:1:0]$ & $(a,b,c,d)=(u,1,1,v)$ \\
$L_4$ & $[1:0:-1:-1]$ & $(a,b,c,d)=(1,u,-1+v,-1)$ \\
$L_5$ & $[0:1:1:1]$ & $(a,b,c,d)=(u,1,1+v,1)$ \\
$L_6$ & $[1:-1:0:0]$ & $(a,b,c,d)=(1,-1,u,v)$ \\
$L_7$ & $[1:2:-1:1]$ & $(a,b,c,d)=(1,2,-1+u+v,1+u)$ \\
\bottomrule
\end{tabular}
\end{center}

\subsection{Corrected table}

\begin{center}
\begin{tabular}{@{}ccccc@{}}
\toprule
Line & PDF type & Exact $\mu$ & Correct type & Generic chain length \\
\midrule
$L_1$ & \wrong{$cA_8$} & 3 & $cA_3$ & 2 \\
$L_2$ & \wrong{$cA_4$} & 3 & $cA_3$ & 2 \\
$L_3$ & \wrong{$cA_7$} & 3 & $cA_3$ & 2 \\
$L_4$ & \wrong{$cA_8$} & 4 & $cA_4$ & 2 \\
$L_5$ & $cA_1$ & 1 & $cA_1$ & 1 \\
$L_6$ & \wrong{$cA_7$} & 5 & $cA_5$ & 3 \\
$L_7$ & $cA_1$ & 1 & $cA_1$ & 1 \\
\bottomrule
\end{tabular}
\end{center}

For the standard local branch model $x^2+y^{n+1}$, a crepant blowup of the double curve
reduces $n$ by two.  Hence the generic chain length is $\lceil n/2\rceil$.  The corrected
generic contribution is therefore
\[
  2+2+2+2+1+3+1=13.
\]

\begin{mistake}
Equation (4.1) and Table 5.1 of the PDF claim 20 old-line-chain blowups based on the
incorrect cA table.  The generic line calculation contributes 13, not 20.  Special
intersection centers must then be recomputed in the corrected order; they cannot be copied
unchanged from the old ledger.
\end{mistake}

\section{The omitted isolated ODP}

Work in the affine chart $d=1$ and translate $P_*$ to the origin:
\[
  a=4+u,\qquad b=-2+v,\qquad c=3+w.
\]
The Hessian matrix of the translated branch germ at the origin is
\[
H=
\begin{pmatrix}
72&1008&720\\
1008&1440&864\\
720&864&1440
\end{pmatrix},
\qquad
\det(H)=-859963392\ne0.
\]
The branch germ is therefore a three-variable node.  Adding the $R^2$ term gives a
nondegenerate quadratic form in four variables, so the point of $Y_0$ above $P_*$ is an
ordinary double point.

The point $P_*$ lies on none of $L_1,\ldots,L_7$ and none of the six listed line-intersection
point centers.  Consequently all divisorial blowups in the PDF are isomorphisms in a
neighborhood of this ODP.

At the point $[a:b:c:d:R]=[4:-2:3:1:0]$, the first base equation of each advertised
small-resolution divisor evaluates as follows:
\begin{center}
\begin{tabular}{@{}lcc@{}}
\toprule
Divisor & First equation & Value at $P_*$ \\
\midrule
$D_{ab,+}$ & $b-a$ & $-6$ \\
$D_{b,-}$ & $b$ & $-2$ \\
$D_{d,-}$ & $d$ & $1$ \\
\bottomrule
\end{tabular}
\end{center}
Thus none of the three divisors contains the ODP above $P_*$.

\begin{mistake}
The Macaulay2 appendix checks only that the three divisor ideals are contained in the double
cover.  It does not check that they contain every node, that their blowups resolve every
node, or that the resulting threefold is smooth.  Divisor containment alone is insufficient.
\end{mistake}

\section{Comparison with the Rust blowup audit}

\subsection{What the 33-stage file actually proves}

The Rust path recorded in \texttt{seven\_line\_stage33\_raw.json} consists of 33 global
curve-center stages and 105 chartwise transforms.  Every recorded center has codimension
two and branch multiplicity two, so every recorded stage satisfies
\[
  2-1-\left\lfloor\frac{2}{2}\right\rfloor=0.
\]
This proves crepantness of the recorded stages, but the file was written with its terminal
audit intentionally disabled.  It is a checkpoint, not a resolution certificate, and its
blowup order is not the 36-center order claimed by the PDF.

On charts 212 and 213, exact reduction proves that two one-dimensional components remain.
For example, on chart 212 one component is
\[
  C_2:\quad x_2+x_3^2=0,\qquad x_0x_3^2-x_3-1=0.
\]
Reducing the branch polynomial and all three partial derivatives modulo this ideal gives
zero.  The corresponding chart-213 equations are
\[
  x_0x_3^2+1=0,\qquad x_2+x_3+1=0.
\]

The four bounded-grid samples previously treated as four points are paired by the overlap:
\begin{align*}
(0,-1,-1)_{212}&=(-1,0,-1)_{213},\\
(2,-1,1)_{212}&=(-1,-2,1)_{213}.
\end{align*}
Both global samples lie on $C_2$ and therefore are not isolated ODPs.

\begin{mistake}
The earlier audit report was wrong to call the stage-33 state a ``final singular locus'' and
wrong to report a global terminal ODP count of zero.  What was proved is narrower: zero of
those four chart samples are ODPs.  The global terminal count is unknown, and the initial
ODP above $P_*$ was omitted from that report.
\end{mistake}

\subsection{Why the two ledgers cannot yet be merged}

The PDF ledger begins with incorrect generic line data and does not include transform-by-
transform equations.  The Rust ledger uses a different curve-first order, omitted the
isolated point, and stops with nonlinear singular curves.  Therefore neither existing ledger
is a complete verified resolution, and their center counts cannot be combined arithmetically.

\section{Status of the downstream claims}

\renewcommand{\arraystretch}{1.18}
\begin{longtable}{@{}p{0.29\textwidth}p{0.22\textwidth}p{0.43\textwidth}@{}}
\toprule
Claim & Status & Reason \\
\midrule
\endfirsthead
\toprule
Claim & Status & Reason \\
\midrule
\endhead
Projection gives $R^2=\Delta$ & Verified & Direct algebraic identity. \\
Seven lines plus $P_*$ & Verified & Exact projective Jacobian decomposition; direct substitution. \\
PDF transverse cA table & \wrong{False} & Five entries disagree with exact local Milnor numbers. \\
20 old line-chain blowups & \wrong{False} & Correct generic contribution is 13. \\
36-center ledger resolves the branch & \wrong{Not established} & The ledger is hardcoded; no successive strict-transform audit is supplied. \\
Only four final nodes remain & \wrong{False as written} & The ODP above $P_*$ is untouched and absent from the node list. \\
Three divisors give a smooth projective model & \wrong{False as written} & None of the divisors contains the ODP above $P_*$. \\
$\#Y_0(\mathbb F_p)$ values & Separately checkable & These count the singular double cover and do not require a resolution. \\
$58p^2+82p$ correction & \wrong{Unsupported} & It depends on the invalidated exceptional-stratum ledger. \\
$h^{1,1}=66$, $e=132$, rigidity & \wrong{Unsupported} & These depend on the same unresolved geometry. \\
Level-13 modularity of the claimed $X$ & \wrong{Unsupported here} & The trace sequence uses the unsupported resolved counts. \\
\bottomrule
\end{longtable}

The arithmetic equalities in the PDF may be internally consistent once one assumes
$58p^2+82p$ and $h^{1,1}=66$.  Internal consistency is not a substitute for proving the
geometric correction formula.

\section{Work required for a valid final theorem}

A corrected proof should supply all of the following.
\begin{enumerate}
  \item Start from the corrected transverse table and recompute the ordered strict transforms
  of all seven lines and their special intersections.
  \item Include the isolated ODP above $P_*$ from the outset.  Supply a global Weil divisor
  containing it and verify that its blowup gives a projective small resolution, or use another
  valid projective crepant construction.
  \item At every divisorial stage, record the actual center ideal, branch multiplicity, branch
  transform, active affine charts, and exact singular locus after the transform.
  \item Resolve or account for nonlinear residual curve components such as the chart-212
  component $C_2$ in the Rust order.
  \item Prove terminal smoothness by an exact Jacobian audit on every final chart.
  \item Only then recompute the exceptional-stratum point-count correction, Picard rank,
  Euler characteristic, Hodge numbers, and Frobenius traces.
\end{enumerate}

Until these steps are completed, the correct conclusion is not that Q1 has four final ODPs,
but that the proposed resolution and its numerical consequences remain unproved.

\section{Reproducibility files}

All files used for this audit are in the \texttt{Q1/} directory.

\begin{description}
  \item[\texttt{README.md}] Directory manifest and commands for rerunning the checks.
  \item[\texttt{q1\_revised\_report.tex/pdf}] This report and its compiled PDF.
  \item[\texttt{q1\_revised\_findings.json}] Machine-readable corrected conclusions and
  explicit unresolved fields for the final singular locus and ODP count.
  \item[\texttt{q1\_expository\_report.pdf}] The newly supplied expected-result report.
  \item[\texttt{q1\_branch\_polynomial.txt}] The supplied octic.
  \item[\texttt{q1\_exact\_checks.sing/txt}] Exact transverse and isolated-point checks.
  \item[\texttt{q1\_stage33\_curve\_checks.sh/txt}] Exact residual-curve reductions.
  \item[\texttt{seven\_line\_initial.json}] Initial Rust singular-component output.  It records
  the seven lines but omits $P_*$, as explained above.
  \item[\texttt{seven\_line\_stage33\_raw.json}] Full 33-stage raw checkpoint.
  \item[\texttt{seven\_line\_crepant\_audit.json}] Earlier merged audit.  Its
  \texttt{resolution\_complete=false} flag remains essential; its terminal wording is
  superseded by this report.
\end{description}

\appendix
\section{Exact command output}

The principal output of \texttt{Singular -q q1\_exact\_checks.sing} is
\begin{verbatim}
L1_MILNOR=3 TYPE=cA3
L2_MILNOR=3 TYPE=cA3
L3_MILNOR=3 TYPE=cA3
L4_MILNOR=4 TYPE=cA4
L5_MILNOR=1 TYPE=cA1
L6_MILNOR=5 TYPE=cA5
L7_MILNOR=1 TYPE=cA1
ISOLATED_POINT_BRANCH_VALUE=0
ISOLATED_POINT_MILNOR=1
ISOLATED_POINT_HESSIAN_DET=-859963392
ISOLATED_POINT_DOUBLE_COVER_TYPE=ODP
Dab_plus_FIRST_GENERATOR_AT_POINT=-6
Db_minus_FIRST_GENERATOR_AT_POINT=-2
Dd_minus_FIRST_GENERATOR_AT_POINT=1
\end{verbatim}

The stage-33 curve check reports dimension one and zero remainders for the branch and all
partials modulo each listed chart ideal.

\end{document}
